% Review
\subsection{Závěr analýzy}

% Review
V této části textu shrnu nejdůležitější výstupy provedené analýzy konkurenčních řešení. Tyto výstupy budou sloužit jako podklad pro sestavení požadavků tohoto projektu.

% Review
\subsubsection{Vzhled a uživatelská zkušenost}

% Review
Většina analyzovaných aplikací se snaží poskytnout uživateli přehledný a snadno použitelný způsob učení. Vzhled uživatelského rozhraní je u většiny aplikací minimalistický. Tomu odpovídá i struktura stránek a samotných kurzů, která bývá velice jednoduchá.

% Review
Aplikace často používají animace a další vizuální efekty, aby udělaly učení a používání aplikace vizuálně poutavější. To může potenciálně pomoci udržet pozornost uživatele při učení. Současné řešení animace a efekty téměř nepoužívá a proto stojí za zvážní jejich použití.

% Review
\subsubsection{Způsob výuky}

% Review
Analyzované aplikace volí různé způsoby výuky. Aplikace Mimo a Sololearn učí uživatele pomocí krátkých cvičení, Codecademy pomocí delších úkolů a Scrimba pomocí vysoce interaktivních lekcí. Každý způsob výuky poskytuje určité výhody a nevýhody a bude třeba zvážit, který způsob je pro uživatele nejvhodnější.

% Review
Při učení všechny aplikace poskytují určitou formu nápovědy pomocí umělé inteligence. Ta například může uživateli vysvětlit danou látku, nebo chybu, kterou ve cvičení udělal. Tato funkce může být pro učení a uživatelskou zkušenost klíčová a bude třeba ji do aplikace zahrnout.

% Review
V rámci učení jsou uživateli zobrazovány různé typy cvičení. Typicky se jedná o doplňování slov do kódu, psaní kódu do editoru a výběr správné odpovědi z nabídky.

% Review
Ve většině aplikací je uživateli látka vysvětlena pomocí krátkých textových popisů. Výjimkou je aplikace Scrimba, ve které je látka uživateli vysvětlena pomocí audio nahrávky v kombinaci s prezentací a interaktivním editorem. Použití audio nahrávek a interaktivních prvků může být pro uživatele potenciálně poutavější a srozumitelnější, než pouhé čtení textu. Podobnou formu výuky bude vhodné zvážit při návrhu struktury učení v této práci.

% Review
Většina aplikací také zobrazuje lekce kurzu ve formě seznamu. To se zásadně liší od současné aplikace, ve které jsou lekce vizuálně uspořádány do tvaru připomínající cestu kurzem. Přestože tato struktura může být vizuálně poutavější, a například aplikace Mimo ji také využívá, bude třeba zvážit, zda-li by pro uživatele nebylo přehlednější uspořádání lekcí do přehlednějšího seznamu.

% Review
\subsubsection{Motivace uživatelů}

% Review
Ve většině aplikací jsou používány určité funkce, které mají za cíl motivovat uživatele, aby se učil a dlouhodobě používal danou aplikaci.

% Review
Mezi tyto prvky se řadí zejména různé způsoby gamifikace. Mezi nejpoužívanější způsoby v analyzovaných aplikacích patří například streak, body XP, ligy učení, srdíčka při učení a virtuální měna. 

% Review
Kromě těchto prvků aplikace také často využívají indikátory pokroku v rámci daných kurzů. Dokončené lekce jsou například označovány zelenou barvou a uživateli je zobrazován procentuální pokrok v daném kurzu. Tyto ukazatele pokroku ve stávající aplikaci nejsou vůbec zahrnuty a stojí za zvážení jejich použití.

% Review
Posledním výrazným prvkem pro motivaci uživatelů jsou certifikáty o dokončení daného kurzu. Po dokončení kurzu a splnění podmínek si může uživatel ve většině aplikací stáhnout certifikát o jeho dokončení ve formátu PDF. Tento certifikát pak může uživatel například využít při hledání práce nebo prezentaci svých znalostí. Tuto funkci bude opět vhodné zvážit při sestavování požadavků tohoto projektu.  

% Review
\subsubsection{Odlišení aplikace od existujících řešení}

% Review
Na závěr analýzy se zaměřím na specifické funkce, kterými se bude moci tato aplikace odlišit od analyzovaných konkurenčních řešení. Tyto funkce budou moci aplikaci poskytnout na trhu konkurenční výhodu a uživatelům lepší uživatelskou zkušenost.

% Review
\paragraph{Interaktivní lekce s audio výkladem}

% Review
Hlavní konkurenční výhodou této aplikace budou interaktivní lekce, které kombinují výhody přístupů analyzovaných aplikací. Lekce budou krátké a budou se skládat z interaktivních stránek výkladu a cvičení. Způsob výuky tak bude podobný způsobu výuky aplikací Mimo a Sololearn. Průběh učení ovšem bude doplněn o audio nahrávku s výkladem, podobně jako tomu je v aplikaci Scrimba. Uživatelům tak bude látka prezentována po malých krocích, což jim umožní látku snadněji pochopit. Aplikaci to dále umožní přesně identifikovat, se kterými koncepty má daný uživatel problém a tyto koncepty mu bude moci častěji ukazovat při opakování. Audionahrávky doplněné o animace, kód a obrázky navíc budou moci uživatelům přinést lepší uživatelskou zkušenost. Uživatelé tak nebudou muset číst dlouhé odstavce textu, ale veškerá látka jim bude tímto způsobem odprezentována a názorně vysvětlena.

% Review
\paragraph{Důraz na opakování látky}

% Review
Další konkurenční výhodou oproti analyzovaným aplikacím bude větší důraz na opakování již naučené látky. Většina konkurenčních aplikací na opakování klade velmi malý důraz. Pro uživatele tak může být obtížné si koncepty a informace dlouhodobě zapamatovat. 

% Review
Stávající aplikace již obsahuje určité formy opakování pomocí metody spaced repetition. % Write - zde dát odkaz na Riskinovu bakalářku
Tuto formu opakování bude třeba změnit a přizpůsobit tak, aby bylo možné ji využít pro výuku softwarového inženýrství. Při návrhu aplikace tak bude třeba zvážit, jakým způsobem bude opakování látky v aplikaci prováděno.

% Review
\paragraph{Zaměření na gamifikaci}

% Review
Některé konkurenční aplikace využívají prvky gamifikace pro zvýšení a udržení motivace uživatele. Tyto prvky nejsou ovšem zdaleka použity tak efektivně, jak by mohly. Některé aplikace tyto prvky využívají velice zřídka a jiné je zase nevyužívají způsobem, jaký by byl pro uživatele optimální.

% Review
Aplikace se tak může na tento aspekt učení zaměřit a poskytnout uživatelům efektivnější a propracovanější použití gamifikace.

% Review
\paragraph{Zaměření na uživatelskou zkušenost}

% Review
Některé aplikace, jako například Sololearn nebo Scrimba, mají řadu nedostatků, co se návrhu vzhledu aplikace a uživatelské zkušenosti týče. Tyto nedostatky mohou být pro uživatele odrazující. Při vývoji této aplikace tak bude možné se více zaměřit na vzhled, uživatelskou zkušenost a testování, díky čemuž by aplikace mohla poskytnout kvalitnější uživatelskou zkušenost.

% Review
\paragraph{Kvalitní učení}

% Review
Efektivní a kvalitní učení je klíčovým aspektem vzdělávacích aplikací. Přesto konkurenční aplikace obsahují řadu nedostatků, které mohou uživateli znepříjemnit učení. Při zapnutí aplikace Sololearn je například uživatel odkázán na svůj profil, kde je zahlcen řadou nepodstatných informací. Místo toho by například mohl být odkázán přímo na stránku kurzu, kde by mohl pokračovat v učení.

% Review
V průběhu učení je například v aplikaci Codecademy zobrazeno uživateli značné množství materiálů a informací, které pro něj mohou být velice rozptylující a snižovat tak efektivitu samotného učení.

% Review
Aplikace by se tak mohla zaměřit na předejití těchto nedostatků a poskytnout uživateli příjemnější a kvalitnější výuku.

% Review
\paragraph{Studijní plány}

% Review
V rámci učení se jednoduše může stát, že se uživatel potřebuje učit několik kurzů zároveň. Začínající programátor se například musí naučit řadu technologií a jazyků, se kterými se musí naučit pracovat. 

% Review
Pokud se ovšem uživatel chce učit více kurzů zároveň, musí mezi jednotlivými kurzy neustále složitě přepínat a sám plánovat, jakou látku se bude učit dál.

% Review
Aplikace by tak mohla poskytnout uživateli možnost si zapsat jednotlivé kurzy a zvolit si úroveň, které chce uživatel v kurzech dosáhnout. Následně by mu mohla předkládat lekce k naučení podle automaticky vygenerovaného studijního plánu. Uživatel by si tak mohl zapsat kurzy a následně se jen učit lekce, které mu aplikace doporučí.

% Review
\paragraph{Sociální aspekty}

% Review
Další slabou stránkou aplikací je nedostatek sociálních interakcí mezi uživateli. Aplikace často neobsahují uživatelské profily, sledování aktivit ostatních uživatelů, vyhledávání jejich profilů a sdílení pokroku mezi ostatními uživateli. Tento sociální aspekt používání může být zásadním zdrojem motivace pro některé uživatele a mohl by přispět ke zlepšení uživatelské zkušenosti.

% Review
\paragraph{Využití animací, přechodů a efektů}

% Review
Konkurenční aplikace zřídka využívají animace, přechody a efekty mezi stránkami. Přestože je třeba používat tyto prvky v přiměřeném množství, může jejich použití docílit toho, že bude aplikace pro uživatele poutavější a bude se tím moci odlišit od konkurentů.
